\documentclass[10pt,a4paper]{amsart}
\usepackage{amsmath,amssymb,amsthm,graphicx}
\usepackage{verbatim}
\newenvironment{reference}{\comment}{\endcomment}
%\newenvironment{reference}{}{}
\newenvironment{slogan}{\comment}{\endcomment}
\newenvironment{history}{\comment}{\endcomment}

% For commutative diagrams we use Xy-pic
\usepackage[all]{xy}

% We use 2cell for 2-commutative diagrams.
\xyoption{2cell}
% Theorem environments.
%
\theoremstyle{plain}
\newtheorem{theorem}[subsection]{Theorem}
\newtheorem{proposition}[subsection]{Proposition}
\newtheorem{lemma}[subsection]{Lemma}

\theoremstyle{definition}
\newtheorem{definition}[subsection]{Definition}
\newtheorem{example}[subsection]{Example}
\newtheorem{exercise}[subsection]{Exercise}
\newtheorem{situation}[subsection]{Situation}

\theoremstyle{remark}
\newtheorem{remark}[subsection]{Remark}
\newtheorem{remarks}[subsection]{Remarks}

\numberwithin{equation}{subsection}

% Macros
%
\def\lim{\mathop{\mathrm{lim}}\nolimits}
\def\colim{\mathop{\mathrm{colim}}\nolimits}
\def\Spec{\mathop{\mathrm{Spec}}}
\def\Hom{\mathop{\mathrm{Hom}}\nolimits}
\def\Ext{\mathop{\mathrm{Ext}}\nolimits}
\def\SheafHom{\mathop{\mathcal{H}\!\mathit{om}}\nolimits}
\def\SheafExt{\mathop{\mathcal{E}\!\mathit{xt}}\nolimits}
\def\Sch{\mathit{Sch}}
\def\Mor{\mathop{\mathrm{Mor}}\nolimits}
\def\Ob{\mathop{\mathrm{Ob}}\nolimits}
\def\Sh{\mathop{\mathit{Sh}}\nolimits}
\def\NL{\mathop{N\!L}\nolimits}
\def\CH{\mathop{\mathrm{CH}}\nolimits}
\def\proetale{{pro\text{-}\acute{e}tale}}
\def\etale{{\acute{e}tale}}
\def\QCoh{\mathit{QCoh}}
\def\Ker{\mathop{\mathrm{Ker}}}
\def\Im{\mathop{\mathrm{Im}}}
\def\Coker{\mathop{\mathrm{Coker}}}
\def\Coim{\mathop{\mathrm{Coim}}}

% Boxtimes
%
\DeclareMathSymbol{\boxtimes}{\mathbin}{AMSa}{"02}

%
% Macros for moduli stacks/spaces
%
\def\QCohstack{\mathcal{QC}\!\mathit{oh}}
\def\Cohstack{\mathcal{C}\!\mathit{oh}}
\def\Spacesstack{\mathcal{S}\!\mathit{paces}}
\def\Quotfunctor{\mathrm{Quot}}
\def\Hilbfunctor{\mathrm{Hilb}}
\def\Curvesstack{\mathcal{C}\!\mathit{urves}}
\def\Polarizedstack{\mathcal{P}\!\mathit{olarized}}
\def\Complexesstack{\mathcal{C}\!\mathit{omplexes}}
% \Pic is the operator that assigns to X its picard group, usage \Pic(X)
% \Picardstack_{X/B} denotes the Picard stack of X over B
% \Picardfunctor_{X/B} denotes the Picard functor of X over B
\def\Pic{\mathop{\mathrm{Pic}}\nolimits}
\def\Picardstack{\mathcal{P}\!\mathit{ic}}
\def\Picardfunctor{\mathrm{Pic}}
\def\Deformationcategory{\mathcal{D}\!\mathit{ef}}
\usepackage{cleveref}
\crefname{originallemma}{original lemma}{original lemmas}
\Crefname{originallemma}{Original lemma}{Original lemmas}
\makeatletter
\newlabel{coherent-lemma-flat-base-change-cohomology}{{02KH}{source}{}{sourcepointer-1}{}}
\newlabel{coherent-lemma-flat-base-change-cohomology@cref}{{[originallemma][0][]02KH}{[1][0][]source}{}{}{}}
\makeatother
\crefname{originallemma}{original lemma}{original lemmas}
\Crefname{originallemma}{Original lemma}{Original lemmas}
\makeatletter
\newlabel{fields-lemma-separable-degree}{{09HJ}{source}{}{sourcepointer-2}{}}
\newlabel{fields-lemma-separable-degree@cref}{{[originallemma][0][]09HJ}{[1][0][]source}{}{}{}}
\makeatother
\crefname{originalreference}{original reference}{original references}
\Crefname{originalreference}{Original reference}{Original references}
\makeatletter
\newlabel{lemma-genus-goes-down}{{0CE3}{source}{}{sourcepointer-3}{}}
\newlabel{lemma-genus-goes-down@cref}{{[originalreference][0][]0CE3}{[1][0][]source}{}{}{}}
\makeatother
\crefname{originallemma}{original lemma}{original lemmas}
\Crefname{originallemma}{Original lemma}{Original lemmas}
\makeatletter
\newlabel{morphisms-lemma-normalization-in-terms-of-components}{{0CDV}{source}{}{sourcepointer-4}{}}
\newlabel{morphisms-lemma-normalization-in-terms-of-components@cref}{{[originallemma][0][]0CDV}{[1][0][]source}{}{}{}}
\makeatother
\crefname{originallemma}{original lemma}{original lemmas}
\Crefname{originallemma}{Original lemma}{Original lemmas}
\makeatletter
\newlabel{varieties-lemma-regular-functions-proper-variety}{{04L2}{source}{}{sourcepointer-5}{}}
\newlabel{varieties-lemma-regular-functions-proper-variety@cref}{{[originallemma][0][]04L2}{[1][0][]source}{}{}{}}
\makeatother
\title{The geometric genus of a proper curve is bounded by the separable degree times the first cohomology dimension}
\author{The Stacks Project Authors}
\date{stacks-project:curves.tex:local-snapshot-20261012; GNU FDL 1.2 or later}
\begin{document}
\maketitle
\noindent\textit{Editorial note.} All delivered mathematical text is the author's own, copied from the pinned Stacks Project chapter file curves.tex; no statement or proof step was rewritten, reordered or removed.\par\medskip
\noindent\textit{Editorial note.} The counted proof is the author's complete argument: the separable-degree fibre decomposition, the flat base change matching, the normalization decomposition, and the genus-goes-down summation are all retained verbatim.\par\medskip
\noindent\textit{Editorial note.} The label lemma-bound-geometric-genus-curve names the counted statement, and the proof environment immediately following it in the source is delivered verbatim as its proof.\par\medskip
\noindent\textit{Editorial note.} The counted statement is carried by the author's own lemma environment in the chapter source.\par\medskip
\noindent\textit{Allowed standard prerequisites (editorial).}\begin{itemize}
\item Global functions and separable degree. Varieties, Lemma varieties-lemma-regular-functions-proper-variety; Fields, Lemma fields-lemma-separable-degree. Makes kappa = H\textasciicircum{}0(C, O\_C) a finite extension of k and identifies [kappa : k]\_s with the number of k-embeddings into the algebraic closure, i.e. the number of geometric fibres.
\item Base change and normalization decomposition. Cohomology of Schemes, Lemma coherent-lemma-flat-base-change-cohomology; Morphisms, Lemma morphisms-lemma-normalization-in-terms-of-components. Preserves dim H\textasciicircum{}1 under the field base change and decomposes the normalization of each geometric fibre into nonsingular projective models of its components.
\item Genus goes down under normalization. Same chapter, Lemma lemma-genus-goes-down. Supplies the final inequality dim H\textasciicircum{}1(C\_t, O) >= dim H\textasciicircum{}1(C\_t\textasciicircum{}nu, O) comparing the curve with its normalization.
\end{itemize}\medskip
\ifdefined\hypertarget\hypertarget{sourcepointer-1}{}\fi
\noindent\textit{Source reference.} Original lemma 02KH: sources/stacks-chow/tags.txt, source lines 3372--3372. Prerequisite pointer invoked at proof line(s) 3943; it names a standard fact of a sibling chapter that the counted argument uses, and every step of the counted proof built on it is retained verbatim in the delivered body.\par\medskip
\ifdefined\hypertarget\hypertarget{sourcepointer-2}{}\fi
\noindent\textit{Source reference.} Original lemma 09HJ: sources/stacks-chow/tags.txt, source lines 12007--12007. Prerequisite pointer invoked at proof line(s) 3921; it names a standard fact of a sibling chapter that the counted argument uses, and every step of the counted proof built on it is retained verbatim in the delivered body.\par\medskip
\ifdefined\hypertarget\hypertarget{sourcepointer-3}{}\fi
\noindent\textit{Source reference.} Original reference 0CE3: sources/stacks-curves/pinned-original-source.tex, source lines 3850--3850. Prerequisite pointer invoked at proof line(s) 3950; it names a standard fact of the same chapter that the counted argument uses, and every step of the counted proof built on it is retained verbatim in the delivered body.\par\medskip
\ifdefined\hypertarget\hypertarget{sourcepointer-4}{}\fi
\noindent\textit{Source reference.} Original lemma 0CDV: sources/stacks-chow/tags.txt, source lines 15555--15555. Prerequisite pointer invoked at proof line(s) 3947; it names a standard fact of a sibling chapter that the counted argument uses, and every step of the counted proof built on it is retained verbatim in the delivered body.\par\medskip
\ifdefined\hypertarget\hypertarget{sourcepointer-5}{}\fi
\noindent\textit{Source reference.} Original lemma 04L2: sources/stacks-chow/tags.txt, source lines 5993--5993. Prerequisite pointer invoked at proof line(s) 3919; it names a standard fact of a sibling chapter that the counted argument uses, and every step of the counted proof built on it is retained verbatim in the delivered body.\par\medskip
\begin{lemma}
\label{lemma-bound-geometric-genus-curve}
Let $k$ be a field. Let $C$ be a proper curve over $k$.
Set $\kappa = H^0(C, \mathcal{O}_C)$. Then
$$
[\kappa : k]_s \dim_\kappa H^1(C, \mathcal{O}_C) \geq g_{geom}(C/k)
$$
\end{lemma}

\begin{proof}
Varieties, Lemma \ref{varieties-lemma-regular-functions-proper-variety}
implies $\kappa$ is a field and a finite extension of $k$.
By Fields, Lemma \ref{fields-lemma-separable-degree}
we have $[\kappa : k]_s = |\Mor_k(\kappa, \overline{k})|$
and hence $\Spec(\kappa \otimes_k \overline{k})$ has
$[\kappa : k]_s$ points each with residue field $\overline{k}$.
Thus
$$
C_{\overline{k}} =
\bigcup\nolimits_{t \in \Spec(\kappa \otimes_k \overline{k})} C_t
$$
(set theoretic union). Here
$C_t = C \times_{\Spec(\kappa), t} \Spec(\overline{k})$ where
we view $t$ as a $k$-algebra map $t : \kappa \to \overline{k}$.
The conclusion is that $g_{geom}(C/k) = \sum_t g_{geom}(C_t/\overline{k})$
and the sum is over an index set of size $[\kappa : k]_s$.
We have
$$
H^0(C_t, \mathcal{O}_{C_t}) = \overline{k}
\quad\text{and}\quad
\dim_{\overline{k}} H^1(C_t, \mathcal{O}_{C_t}) =
\dim_\kappa H^1(C, \mathcal{O}_C)
$$
by cohomology and base change
(Cohomology of Schemes, Lemma \ref{coherent-lemma-flat-base-change-cohomology}).
Observe that the normalization $C_t^\nu$ is the disjoint
union of the nonsingular projective models of the
irreducible components of $C_t$
(Morphisms, Lemma \ref{morphisms-lemma-normalization-in-terms-of-components}).
Hence $\dim_{\overline{k}} H^1(C_t^\nu, \mathcal{O}_{C_t^\nu})$
is equal to $g_{geom}(C_t/\overline{k})$.
By Lemma \ref{lemma-genus-goes-down} we have
$$
\dim_{\overline{k}} H^1(C_t, \mathcal{O}_{C_t}) \geq
\dim_{\overline{k}} H^1(C_t^\nu, \mathcal{O}_{C_t^\nu})
$$
and this finishes the proof.
\end{proof}

\end{document}
